\documentclass[12pt]{article}
\usepackage[T1]{fontenc}
\usepackage[utf8]{inputenc}
\usepackage{lmodern}
\usepackage{textcomp}
\usepackage{enumitem}
\usepackage{geometry}
\usepackage{graphicx}
\usepackage{amsmath, amssymb}
\geometry{margin=1in}
\begin{document}
\begin{center}
Chemistry - Graduate Level Assessment 10 \
Total Marks: 40 \
Answer all questions.
\end{center}

\section*{Multiple Choice (10 marks)}
\begin{enumerate}[label=\arabic*.]
\item Calculate the normality of a solution containing 2.45 g of sulfuric acid in 2.00 liters of solution. (MW of H\_2 SO\_4 = 98.1.)
\begin{enumerate}[label=(\alph*)]
    \item 0.0125 equiv/liter
    \item 0.060 equiv/liter
    \item 0.050 equiv/liter
    \item 0.100 equiv/liter
    \item 0.125 equiv/liter
\end{enumerate}
\item Calculate the equilibrium polarization of 13 C nuclei in a 20.0 T magnetic field at 300 K.
\begin{enumerate}[label=(\alph*)]
    \item 7.39 x 10^-5
    \item 9.12 x 10^-5
    \item 3.43 x 10^-5
    \item 6.28 x 10^-5
    \item 1.71 x 10^-5
\end{enumerate}
\item $ \text { If we locate an electron to within } 20 \mathrm{pm} \text {, then what is the uncertainty in its speed? } $\$
\begin{enumerate}[label=(\alph*)]
    \item 8.1$10^6 \mathrm{~m} \cdot \mathrm{s}^{-1}$
    \item 9.4$10^7 \mathrm{~m} \cdot \mathrm{s}^{-1}$
    \item 6.2$10^7 \mathrm{~m} \cdot \mathrm{s}^{-1}$
    \item 1.1$10^8 \mathrm{~m} \cdot \mathrm{s}^{-1}$
    \item 3.7$10^7 \mathrm{~m} \cdot \mathrm{s}^{-1}$
\end{enumerate}
\item Compare the bond order of He\_2 and He\_2^+.
\begin{enumerate}[label=(\alph*)]
    \item Bond order of He\_2 is 0 and He\_2^+ is 0.5
    \item Bond order of He\_2 is 0.5 and He\_2^+ is 0
    \item Bond order of He\_2 and He\_2^+ is 0.5
    \item Bond order of He\_2 is 2 and He\_2^+ is 1.5
    \item Bond order of He\_2 is 1 and He\_2^+ is 0
\end{enumerate}
\item A solution of sulfurous acid, H$_{2}$SO$_{3}$, is present in an aqueous solution. Which of the following represents the concentrations of three different ions in solution?
\begin{enumerate}[label=(\alph*)]
    \item [H$_{2}$SO$_{3}$] \textgreater{} [HSO 3-] = [SO 32-]
    \item [HSO 3-] \textgreater{} [SO 32-] \textgreater{} [H$_{2}$SO$_{3}$]
    \item [SO 32-] \textgreater{} [H$_{2}$SO$_{3}$] \textgreater{} [HSO 3-]
    \item [HSO 3-] = [SO 32-] \textgreater{} [H$_{2}$SO$_{3}$]
    \item [SO 32-] = [HSO 3-] \textgreater{} [H$_{2}$SO$_{3}$]
\end{enumerate}
\end{enumerate}

\section*{True / False (10 marks)}
\textit{Answer True or False}
\begin{enumerate}[label=\arabic*.]
\setcounter{enumi}{5}
\item True or False: The volume of a cube with edges measuring 150.0 mm is 3375 cubic centimeters.
\item True or False: The molecule NO 2- exhibits sp 2 hybridization around its central nitrogen atom.
\item True or False: The recoil energy of an atom with a mass of 10 following the emission of a 2 MeV gamma ray is 214 eV.
\item True or False: During the reduction of NO(g), NO 2(g) is produced without any nitrogen compounds being oxidized.
\item True or False: No two electrons in an atom can have the same set of four quantum numbers due to the Pauli exclusion principle.
\end{enumerate}

\section*{Long Form Response (20 marks)}
\textit{Answer each question with a clear and complete explanation.}
\begin{enumerate}[label=\arabic*.]
\setcounter{enumi}{10}
\item Discuss how the standard molar enthalpy of sulfur dioxide gas changes when the temperature is increased from 298.15 K to 1500 K, and calculate the increase in enthalpy.
\item Discuss the effects of increasing H$_{3}$O+ concentration and the subsequent saturation of H$_{2}$S in a solution containing $Zn^2$+ and $Cu^2$+ ions on the precipitation of sulfides, specifically CuS and ZnS.
\end{enumerate}

\vfill
\noindent\textit{End of Paper}
\end{document}